\noindent\hspace*{\parindent}Zum Verständnis der Positionen (3) und (8) dient Folgendes. Das älteste der jetzigen Mitglieder trägt die Nummer 135: Die ersten 134 sind sämtlich verstorben oder auf andere Weise ausgeschieden. Für jedes dieser 134 ausgeschiedenen Mitglieder ist berechnet worden: der Betrag der von ihm geleisteten Beiträge und, wo Hinterbliebene Pensionen erhalten haben, der Betrag der Witwenpensionen und der Waisenpensionen. Alle diese Zahlungen sind jedoch nicht mit ihren tatsächlichen Beträgen angesetzt worden, sondern nach den gegenwärtig bestehenden Sätzen, nämlich 15 Thaler als jährlicher Beitrag und 200 Thaler als jetziger Betrag der festen Pension. Außerdem sind sämtliche gezahlten oder empfangenen Gelder durch Diskontierung auf den Zeitpunkt des Eintritts des betreffenden Mitglieds zurückgeführt worden. Hinsichtlich der Änderungen, die an diesen Rechnungen noch bei denjenigen der 134 Mitglieder vorgenommen werden mussten, deren Witwen noch leben, verweise ich auf meine Denkschrift von 1845.

So verstanden beträgt für alle diese 134 Mitglieder
\begin{center}
\begin{tabular}{l r l}
die Summe aller Beiträge . . . . . .&25506&Thaler\\
die Summe aller Witwenpensionen&66771&,,\\
die Summe aller Waisenpensionen&9604&,,
\end{tabular}
\end{center}
\noindent woraus sich folgende Durchschnittswerte ergeben: für die Beiträge 190,35 Thaler, für die Witwenpensionen 498,29 Thaler, für die Waisenpensionen 71,67 Thaler. Im Jahr 1845, als die bis dahin ausgeschiedenen Mitglieder in ununterbrochener Folge nur bis zu Nr. 108 reichten, hatte ich für sie eine ganz ähnliche Rechnung ausgeführt. Für die Durchschnittswerte hatte sie, in dieselbe Form gebracht, ergeben: 193,18 Thaler; 503,08 Thaler; 81,07 Thaler. Man erkennt daraus mit Befriedigung, dass sich die bisherigen Erfahrungen schon sehr gut zur Feststellung von Durchschnittswerten eignen und damit eine wertvolle Grundlage zur Abschätzung des künftigen Bedarfs geben. Nach dem neuen Ergebnis ist das Verhältnis der Waisenpension zur Witwenpension im Durchschnitt sehr nahe \(\tfrac{1}{7}\); dies ist zur Feststellung der Position (2) benutzt worden (man vergleiche den Schluss von S. [165]).

Den allmählichen Leistungen jedes neu eintretenden Mitglieds einerseits und den Pensionsbezügen seiner Hinterbliebenen andererseits sind nach den Durchschnittswerten die Summen 190,35 und 569,96 Thaler gleichwertig, die zum Zeitpunkt des Eintritts einmal in die Kasse eingezahlt beziehungsweise aus ihr ausgezahlt werden. Wenn jedes Jahr ein neues Mitglied einträte, brauchte man sich daher nur diese beiden Zahlungen jedes Jahr wiederholt zu denken und sie dann zu kapitalisieren, das heißt, der Kasse ein für alle Mal die Kapitalsumme 5439 Thaler als Einnahme zuzuweisen und die anderen 16285 Thaler als Ausgabe anzurechnen, um den Verpflichtungen und Berechtigungen aller künftig beitretenden Mitglieder Rechnung zu tragen. Diese Summen müssen aber noch mit derjenigen Zahl multipliziert werden, die die Durch%
\SourceLeaf{172}schnittszahl der jährlich neu beitretenden Mitglieder ausdrückt. In der Denkschrift von 1845 habe ich dafür \(2\tfrac{1}{2}\) angenommen, ohne zu verschweigen, dass diese Größe sehr ungewiss ist. Alles wohl erwogen, habe ich geglaubt, auch diesmal dieselbe Zahl beibehalten zu müssen; so haben sich die in der Bilanzrechnung angesetzten Positionen (8) und (3) ergeben.
